% Fragmento de entrega (sin preambulo; se compila desde main.tex).
\section*{Bitácora de decisiones -- Problema 2}

Cada decisión registra qué se decidió, qué se descartó, por qué y con qué
evidencia. Todas fueron redactadas por el asistente con evidencia propia y
aceptadas por el alumno; quedan como decisiones vigentes del trabajo.

\subsection*{Decisión 1 -- Esquema de validación}

\textbf{Qué decidí:} ventana expansiva con 4 particiones sobre la muestra
ordenada por (dteday, hr), con bloques de 500 horas: entrena 0:1000 y valida
1000:1500; entrena 0:1500 y valida 1500:2000; entrena 0:2000 y valida
2000:2500; entrena 0:2500 y valida 2500:3000, entrenando siempre con horas
anteriores a las de validación.

\textbf{Qué descarté:} un split aleatorio 80/20 (mezcla futuro en el
entrenamiento y no refleja el uso real) y una sola partición temporal (solo
500 horas de validación, alta varianza).

\textbf{Por qué:} en demanda temporal solo es honesto validar hacia el
futuro; con 4 particiones se tienen 2000 horas de validación en momentos
distintos y el entrenamiento crece (1000, 1500, 2000, 2500), reflejando cómo
mejora el modelo con más datos. Evidencia: las fechas impresas en
\path{02_decisiones_p2} (p. ej., P1: entrena 2011-01-01 a 2011-07-31 y
valida 2011-07-31 a 2011-11-19) y \path{resultados/p2_02_particiones.csv}.

\textbf{Qué pierdo:} los bloques cortan dentro de un día, las particiones no
son independientes (entrenamientos solapados) y no se evalúa extrapolación
estacional completa porque la muestra termina el 2012-09-30.

\subsection*{Decisión 2 -- Combinación de kernels}

\textbf{Qué decidí:} la combinación \textbf{(iii) $(k_{hr}\cdot k_{wd})+k_{rbf}$},
con von Mises para hora y día (kappas en rejilla) y RBF para las demás
variables estandarizadas.

\textbf{Qué descarté:} (i) la suma pura $k_{hr}+k_{wd}+k_{rbf}$ (fuerza la
misma forma horaria todos los días) y (ii) el producto total
$k_{hr}\cdot k_{wd}\cdot k_{rbf}$ (obliga también al clima a interactuar con
hora y día).

\textbf{Por qué:} las gráficas de registradas medias por hora separando
\texttt{workingday} muestran formas distintas (doble pico laboral frente a
máximo al mediodía en fin de semana): hay interacción hora $\times$ día. La
validación lo confirma: RMSLE medio (iii) $0.556$ frente a (ii) $0.623$ y
(i) $0.720$. Evidencia: \path{resultados/p2_02_registrado_por_hr.png},
\path{p2_02_registrado_por_weekday.png} y
\path{p2_02_comparacion_kernels.csv}.

\textbf{Qué pierdo:} (iii) asume el efecto del clima aditivo sobre la
interacción hora-día; una interacción clima $\times$ hora por bandas podría
capturar más, a costa de más parámetros.

\subsection*{Decisión 3 -- Objetivo logarítmico y métrica RMSLE}

\textbf{Qué decidí:} trabajar con \textbf{$\log 1\mathrm{p}(\texttt{registered})$}
y elegir modelos por \textbf{RMSLE} (RMSE y MAE se reportan en 04 como
contexto).

\textbf{Qué descarté:} el objetivo crudo con RMSE global.

\textbf{Por qué:} errar 20 bicicletas a las 3:00 (nivel $\approx
18$) no cuesta lo mismo que a las 18:00 (nivel $\approx 188$): el error
absoluto crece con el nivel (RMSE crudo de $10.5$ en madrugada a $58.2$ en la
tarde), mientras que RMSLE es mucho más estable entre franjas y en
validación el modelo en log obtiene RMSLE global $0.404$ frente a $0.570$ del
modelo en crudo. Evidencia: \path{resultados/p2_02_error_por_franja.csv}.

\textbf{Qué pierdo:} RMSLE subestima errores absolutos grandes en horas pico,
y al retransformar con $\operatorname{expm1}$ hay sesgo de retransformación
(la media del log no es el log de la media).

\subsection*{Decisión 4 -- Kernel circular válido elegido}

\textbf{Qué decidí:} \textbf{von Mises} $\exp(\kappa(\cos d-1))$ para
$hr$ y para $weekday$, con kappas independientes en rejilla $\{1,3\}$ (en 03
ganaron $\kappa_{hr}=3$, $\kappa_{wd}=1$).

\textbf{Qué descarté:} $\exp(-d/\ell)$, $\max(0,1-d/2)$ y
$\pi^2/6-\pi d/2+d^2/4$ (los cuatro son válidos según la verificación, pero
no admiten un parámetro de concentración tan interpretable como $\kappa$ y
son menos flexibles para picos horarios).

\textbf{Por qué:} es el kernel circular estándar, $\kappa$ controla el ancho
de forma interpretable y es el que justifica la teoría (Taylor de la
exponencial). Evidencia: \path{resultados/p2_01_coeficientes.csv} (los
cuatro sin coeficientes negativos) y \path{p2_01_eigenvalores.csv}.

\textbf{Qué pierdo:} un único $\kappa$ por dimensión no captura formas
asimétricas intra-día (subida rápida, bajada lenta).

\subsection*{Decisión 5 -- workingday, weathersit y yr}

\textbf{Qué decidí:} sacar \textbf{\texttt{workingday}} de la parte RBF y
conservar \textbf{\texttt{weathersit}} (ordinal numérico) y
\textbf{\texttt{yr}} (binaria) en la parte RBF, junto a season, mnth,
holiday, temp, atemp, hum y windspeed.

\textbf{Qué descarté:} incluir \texttt{workingday} como variable extra.

\textbf{Por qué:} \texttt{workingday} vale 1 exactamente si $weekday\in
1,\dots,5$ y $\texttt{holiday}=0$ (verificado en las 3000/3000 filas): es
función exacta de \texttt{weekday} y \texttt{holiday}, que ya entran por el
kernel de periodo 7 y por la parte RBF. \texttt{weathersit} discrimina el
nivel (161.9, 141.6, 91.8, 22.0) y \texttt{yr} captura el crecimiento
interanual (117.5 a 194.5). Evidencia: sección 4 de
\path{02_decisiones_p2}.

\textbf{Qué pierdo:} \texttt{holiday} queda como binaria en RBF, así que el
efecto de festivos puntuales (Navidad) se captura solo parcialmente.

\subsection*{Decisión 6 -- Estandarización de "otras" y entrada del HGB}

\textbf{Qué decidí:} estandarizar las "otras" \textbf{una sola vez con la
muestra completa de 3000 horas} (para que la matriz M4 de 3000$\times$3000 se
calcule una vez y se rebane) y dar al HGB \textbf{todas las columnas de
entrada} (incluida \texttt{workingday}), porque es la referencia sin
restricción de kernels.

\textbf{Qué descarté:} estandarizar por partición (impide la matriz única
precalculada que pide el enunciado para M4) y limitar el HGB a las variables
de los SVR.

\textbf{Por qué:} la fuga por estandarizar solo covariables (sin el objetivo)
es mínima y despreciable frente a las diferencias entre modelos (RMSLE
$0.44$ frente a $0.60$--$1.00$). El HGB, sin kernels, se beneficia de ver
\texttt{workingday} explícito. Evidencia: secciones 1 de
\path{03_modelos_p2} y 04.

\textbf{Qué pierdo:} un ligero optimismo en validación por la
estandarización global, del orden del error de estimación con 1000--2500
puntos.

\subsection*{Nota sobre kernels no válidos (evidencia, no decisión)}

$K_I$ y $K_{II}$ no son kernels: $a_4\approx-3.19\times10^{-4}$ y
$a_3\approx-0.0712$, con $\lambda_{\min}(24\mathrm{h})\approx-0.00505$ y
$-0.839$. M4 con ellos (03, sección 5) da $\lambda_{\min}=-0.247$ y
$-45.4$ en la matriz de entrenamiento, sin advertencias de libsvm, y RMSLE
de validación $0.552$ y $25.45$ (con todos los puntos como vectores de
soporte en el caso $k_{II}$) frente a $0.439$ del M4 válido: el dual pierde
convexidad, óptimo global y garantía de convergencia de SMO.
